\section{Provenance and Error Analysis}
\label{sec:analysis}

The aggregate result is neutral; the question is what the 38 changes consist of.

\subsection{Corrections are rule-consistent, but the deciding factor is \clean{} confidence}
\label{sec:corrections}

Of the 23 corrections, 16 are \hate{} recoveries (7 from \off{}, 9 from \clean{}) and
7 are \clean{} de-flagging cases. All 16 recoveries have a rule firing with a \hate{}
score above the \off{} score, and the rule \hate{} score has median 76.52\% (mean
75.07\%, range 23.51\%--99.39\%): the evidence is strong on average, though individual scores range down to 23.51\%.

All 7 \clean{} de-flagging cases have a silent rule ($r_H=r_O=0$), so $f_H=0.2\,p_H$
and $p_H\ge0.70$ is needed merely to reach $\tau_H=0.14$; a flagged \hate{} at
$p_H\approx0.5$ ends at $f_H\approx0.10$ and falls back to \clean{}, whose channel keeps
full neural weight. Silence acts only as a multiplicative prior on \hate{}. The
alternative convention (silent rule contributes nothing) changes 30 decisions and
raises macro-F1 to 66.14\% ($p=0.44$; McNemar $p=0.54$); it is undecided, so we keep
the shipped convention for its net corrections ($+8$ vs.\ $+4$).

\subsection{Regressions are rule-driven and directional}
\label{sec:regressions}

The 15 regressions split into 10 rule-driven and 5 threshold-driven cases, aligning
perfectly with direction (Fig.~\ref{fig:results}b). In all 10 rule-driven cases (four
\off{}$\to$\hate{}, six \clean{}$\to$\hate{}) the rule fires towards \hate{}, so fusion
trusts lexical evidence over the neural label and the lexicon is wrong about
intensity or target. The \off{} cases are the expected failure of a phrase lexicon:
profanity used as an intensifier, not an attack. The \clean{} cases are false
positives on non-abusive text; one is a political comment whose charged multi-word
entry fires at $r_H=0.62$ while the transformer prefers \clean{} ($p_C=51.50\%$ vs.\
$p_H=47.30\%$), but the \clean{} channel misses $\tau_C=0.60$ and the rule lifts $f_H$
to $58.80\%$, far above $\tau_H=0.14$, so the decision goes to \hate{}. A lexicon
entry's \emph{political} sensitivity can thus leak into moderation decisions
(Sect.~\ref{sec:discussion}).

The 5 threshold-driven regressions are the mirror image: the rule is silent, the
fused \hate{} score is $0.2\,p_H$, and below $p_H\approx0.7$ the \hate{} channel no
longer wins; two miss by exactly 0.1\% ($f_H=13.90\%$ vs.\ $\tau_H=0.14$). One
mechanism ($\alpha_H=0.2$, $\tau_H=0.14$) thus produces 7 correct and 5 incorrect
de-flagging decisions and cannot be removed without also removing the corrections.
The decomposition is exhaustive---26 of 38 changes rule-driven, 12
threshold-driven. The interpretability of the hybrid model does not rely on heatmaps; instead, it stems from the ability to explicitly identify which channel influenced the result, the direction of that influence, the specific text segments involved, and the magnitude of the impact. This applies only to the 16\% of comments where the rule fires; the remaining 84\% have no lexical span.

\subsection{The corrections fusion does not make}
\label{sec:missed}

When examining aggregate metrics, instances where the rule engine detects no signal are indistinguishable from instances where a signal is detected but subsequently disregarded by the fusion mechanism. In our results, the latter occurred: the rule engine detected a signal, but the Transformer model held a high-confidence \clean{} classification---meaning the rule engine\textquotesingle s signal was insufficient to alter the final outcome. Among the 827 AI errors, 148 non-\clean{}
comments have the rule firing in the correct direction and 151 \clean{} comments have
it correctly silent, so the rule agrees with gold on 299 errors and fusion acts on 23
of them (7.7\%). The remaining \textbf{132 missed corrections} are non-\clean{} cases
where the rule already points right yet the hybrid is still wrong. Separately, 144 of
the 151 silent-\clean{} errors stay wrong because the AI's \hate{} probability is high
enough that even one fifth of it clears $\tau_H$.

The 132 are not missed for weak evidence: their median rule \hate{} score is 76.09\%,
indistinguishable from the 76.52\% median of corrected cases (minimum 23.51\%). The deciding factor is therefore not evidence strength but the \clean{} channel\textquotesingle s confidence. They are missed because
the \clean{} channel keeps full AI weight, so a confidently wrong \clean{} cannot be
outvoted (median fused \clean{} score 97\% for missed vs.\ 46\% for corrected).
Evidence-gated fusion (Sect.~\ref{sec:implications}) would abstain, but we did not tune
such a rule on test.

\subsection{What the rule channel responds to}
\label{sec:cue-taxonomy}

The 16 \hate{} recoveries carry a multi-label cue taxonomy: ten teencode/abbreviated
profanity, five standard-spelling profanity or animal slur, eight derogatory wording
without profanity; categories overlap. The channel is not merely a profanity detector.

Which components do the work? Only
one switch matters: removing per-entry confidence costs 3.4\% standalone macro-F1
(52.58\% to 49.14\%), almost all on \off{}; no other switch moves the engine by more
than 0.8\%. Rescaling severity rewrites 1{,}046 rule scores yet changes no hybrid
decision ($p=1.0$). Dropping obfuscation views trades recall for precision (firings
1{,}052 $\to$ 779, precision 58.75\% $\to$ 69.58\%). Inside the hybrid no switch leaves noise ($|\Delta|\le0.12\%$).

Table~\ref{tab:cases} shows five decisions. Case~1: two exact matches carry \hate{}
from $p_H{=}0.19$ to $f_H{=}0.81$, past the 0.78 \clean{} channel. Case~2: \hate{} 0.01
\emph{below} \off{}, but \off{} misses $\tau_O{=}0.66$ while \hate{} clears $\tau_H{=}0.14$.
Case~3: silent rule, $f_H{=}0.2\,p_H$, fallback \clean{}. Cases~4--5: over-firing on
profanity-as-interjection, and \hate{} missing $\tau_H$ by 0.1\%.

\begin{table}[ht]
\centering
\scriptsize
\setlength{\tabcolsep}{3pt}
\renewcommand{\arraystretch}{1.0}
\caption{Case studies from the 38 changed decisions (\texttt{*} masks profanity).
$H/O/C$ are the fused channel scores, $\tau_c$ the frozen per-channel thresholds and
$r_H$ the aggregated \hate{} rule score; the figure after each span is its per-match
score. Rows 1--3 are corrections, 4--5 regressions.}
\label{tab:cases}
\begin{tabularx}{\textwidth}{@{}l>{\raggedright\arraybackslash}p{6.3cm}>{\raggedright\arraybackslash}X@{}}
\toprule
\# & Comment and rule evidence & Gold $\to$ AI $\to$ Hybrid \\
\midrule
1 & \emph{để avata bít ngay là ngu hơn chó} \newline $r_H{=}0.97$; ngu hơn 0.86, chó 0.79 (exact)
  & \hate{} $\to$ \clean{} $\to$ \hate{} \newline $H\,0.81 > C\,0.78$ \\
2 & \emph{\dots\ cũng đ*o thể yêu thương nổi} \newline $r_H{=}0.71$; cũng đ*o 0.71 (exact)
  & \hate{} $\to$ \off{} $\to$ \hate{} \newline $H\,0.49 < O\,0.50$ \\
3 & \emph{hóng bro tét đ*t thằng tùng} \newline $r_H{=}0$; accent-less alias, weight 0
  & \clean{} $\to$ \hate{} $\to$ \clean{} \newline $H\,0.10 < \tau_H{=}0.14$ \\
\midrule
4 & \emph{\dots\ nói không.. đợi cái loz} \newline $r_H{=}0.87$; cái loz 0.87 (exact)
  & \off{} $\to$ \off{} $\to$ \hate{} \newline $H\,0.70 > O\,0.68$ \\
5 & \emph{\dots abe tuổi lol mà so đc} \newline $r_H{=}0$; no scoring match
  & \hate{} $\to$ \hate{} $\to$ \clean{} \newline $H\,13.90\% < \tau_H{=}0.14$ \\
\bottomrule
\end{tabularx}
\end{table}

